% University of Helsinki
% Department of Mathematics and Statistics
% Template for Bachelor's or Master's thesis
%
% This template is intended as a starting point for a thesis;
% feel free to customize it or use your own template.
% See the comments for hints and customization.
%
% Initial version: 29 August 2024 / Petri Laarne


% Add the 'twoside, openright' options when producing a proper printed version
\documentclass[a4paper, 11pt]{report}

\usepackage[hidelinks]{hyperref} % Always load this: clickable links and PDF reader navigation bar
\usepackage{lastpage} % Used to determine the total page count
\usepackage{calc} % For layout calculations involving TeX length units (template title page)
\usepackage{graphicx} % For embedding image files (JPEG, PNG, PDF)
\usepackage{subcaption} % For subfigures
\usepackage{tikz} % For drawing diagrams and other complex vector graphics

% If you use the LuaLaTeX or XeLaTeX compiler, comment out the following line.
% These compilers offer better support for Unicode, so you should use them if possible.
\usepackage[T1]{fontenc}
\usepackage[most]{tcolorbox}
%\DeclareAutoCiteCommand{textcite}{\textcite}{\textcites}
%\ExecuteBibliographyOptions{autocite=textcite}
\DeclareMathOperator*{\argmax}{arg\,max}

\newtcolorbox[auto counter]{proposition}[1][]{
  enhanced,
  breakable,
  fonttitle=\scshape,
  title={Proposition \thetcbcounter},
  #1
}

% Pass the language of the thesis as the LAST option in the list below.
% This changes the language used for dates and words like 'Chapter' or 'Bibliography'.
% You can remove the languages that do not appear in your work.
\usepackage[finnish,swedish,english]{babel}

\usepackage{csquotes} % This package is suggested by 'babel'

% Bibliography. Change the latter command if you modify the .bib file name.
\usepackage{biblatex}
\addbibresource{Resources/library.bib}

\usepackage{mathtools} % Extended mathematics support (this package extends amsmath)
\usepackage{amssymb} % More mathematical symbols
\usepackage{amsfonts} % Blackboard bold mathematics font

% Theorem environments
\usepackage{amsthm}
\theoremstyle{plain}
\newtheorem{theorem}{Theorem}[section] % Theorems are numbered within chapter and section
\newtheorem{lemma}[theorem]{Lemma} % Lemmas etc. follow the same numbering as theorems
\theoremstyle{definition}
\newtheorem{definition}[theorem]{Definition}
\theoremstyle{remark}
\newtheorem{remark}[theorem]{Remark}

% Some examples of custom mathematical commands
\newcommand{\ZZ}{\mathbb Z} % Blackboard bold Z
\newcommand{\RR}{\mathbb R} % Blackboard bold Z
\newcommand{\dx}{\mathop{}\!\mathrm d x} % Nice spacing (and upright d) for dx in integrals

% Finnish notation for evaluating integrals
% Instead of \int_a^b syntax you must use it as \sij{a}{b}
% (To enable the sub-/superscript syntax would require some more complicated TeX wizardry.)
\newcommand{\sij}[2]{\mathop{\bigg/}\nolimits_{\hspace{-0.65em}#1}^{\hspace{0.05em}#2}\hspace{-0.2em}}


\begin{document}

%
% TITLE PAGE
%
\begin{titlepage}

\noindent
\begin{minipage}[b]{\textwidth-45pt}
University of Helsinki\\
Faculty of Science\\
Bachelors Program in Mathematics
% Helsingin yliopisto\\
% Matemaattis-luonnontieteellinen tiedekunta\\
% Helsingfors universitet\\
% Matematisk-naturvetenskapliga fakulteten\\
\end{minipage}
\hfill\includegraphics[width=45pt]{HYlogo.pdf}

\vspace{4pt}\hrule\vfill

\begin{center}
Bachelor's thesis\\[8pt]
% Kandidaatintutkielma / Maisterintutkielma\\[8pt]
% Kandidatavhandlingen / Magisteravhandlingen\\[8pt]

{\huge\bfseries Harmonic Extension at the Boundary of Half-Space}\\[8pt]

Santeri Väätäjä
\end{center}

\vfill\hrule\vspace{4pt}

\noindent
Supervisor:\hfill June 2025\\
Petri Ola\\
\end{titlepage}

%
% ABSTRACT
%
% This command starts a new page.
% In a print version ('twoside' option passed to \documentclass), the page is on the right side
\pagenumbering{gobble}
\cleardoublepage

\noindent\textbf{Otsikko:} Harmoninen laajennus puoliavaruuteen\\ % Otsikko / Titel
\textbf{Kirjoittaja:} Santeri Väätäjä\\ % Tekijä / Författare
\textbf{Kuukausi ja vuosi:} Kesäkuu 2025\\ % Aika / Datum
\textbf{Sivumäärä:} \pageref*{LastPage}~s.\\[1em] % Sivumäärä: XX~s. / Sidoantal: XX~s.
% The page count is updated automatically, but it might require two compilations to be correct.
% If you customize the backmatter, this might become incorrect - check and correct manually if needed.

% If you include an abstract in a different language, wrap it in a foreignlanguage block
\noindent\textbf{Tiivistelmä:}\\ % Tiivistelmä / Referat
Tutkielma käsittelee Dirichlet ongelmaa, jossa tavoitteena on laajentaa reuna-arvot määrittelevä funktio $n$-ulotteisesta Euklidisesta avaruudesta $n+1$-ulotteiseen puoliavaruuteen siten, että laajennus on harmoninen ja siten ratkaisee Laplacen yhtälön. Tätä kohti tutkielmassa edetään esittelmällä Fourier muunnos ja käyttämällä sitä avuksi Poisson ytimen johtamisessa. Poisson ydintä käytetään Poisson integraalin määrittämiseen. Tämä muodostaa eksplisiittisen muodon Dirichlet ongelman ratkaisulle, kun ratkottava osittaisdifferentaaliyhtälö on Laplacen yhtälö.

Tutkielmassa on kaksi päätulosta. Näistä ensimmäinen todistaa, että Poisson integraali toteuttaa reuna-arvo ongelman kun reuna-arvot määräytyvät funktiosta, joka kuuluu $L^p$-avaruuteen. Tätä tulosta täydennetään lisätuloksella, jonka avulla voidaan näyttää, että ratkaisun suppeneminen noudattaa vahvempaa suppenemiskäsitettä silloin, kun reuna-arvot määrittävä funktio kuuluu tiukemmat säännöllisyysvaatimukset toteuttavaan funktiavaruuteen. Toinen tutkielman päätuloksista laajentaa ensimmäistä tulosta vastaavan tuloksen koskemaan myös ääreellisiä Borel mittoja. Tämä mahdollista monipuolisemman reuna-arvo datan käyttämisen reuna-arvo ongelmassa.

Nämä tulokset antavat eksplisiittisen integrallimuodon funktiolle, joka ratkaisee halutun Dirichlet ongelman. Tulokset valottaa myös sitä miten harmoniset laajennukset käyttäytyvät, kun reuna-arvot määrittävän funktion säännöllisyydelle oletetaan erilaisia ehtoja. Lisäksi tämän kaltaisilla tuloksilla on sovelluksia monissa erilaisissa matematiikan ja fysiikan sovelluksissa. Tuloksia voidaan käyttää esimerkiksi osittaisdifferentaaliyhtälöiden analyysissä ja potentiaaliteoriassa. Työn lopussa käyn läpi vielä joitakin mahdollisia yleistyksiä, joihin tuloksia on mahdollista laajentaa.

\vfill

\noindent\textbf{Avainsanat:} Harmoninen analyysi, Osittaisdifferentaaliyhtälöt, Fourier analyysi % Avainsanat / Nyckelord

\clearpage
\noindent\textbf{Title:} Harmonic Extension at the Boundary of Half-Space\\ % Otsikko / Titel
\textbf{Author:} Santeri Väätäjä\\ % Tekijä / Författare
\textbf{Month and year:} June 2025\\ % Aika / Datum
\textbf{Page count:} \pageref*{LastPage}~pp.\\[1em] % Sivumäärä: XX~s. / Sidoantal: XX~s.
% The page count is updated automatically, but it might require two compilations to be correct.
% If you customize the backmatter, this might become incorrect - check and correct manually if needed.

% If you include an abstract in a different language, wrap it in a foreignlanguage block
\begin{foreignlanguage}{english}
\noindent\textbf{Abstract:}\\ % Tiivistelmä / Referat
In this thesis, I will introduce results for the Dirichlet problem in the Euclidean half-space. The goal is to extend a boundary value function from $n$-dimensional Euclidean space to the whole of the $n+1$-dimensional half-space, so that it solves the Laplace equation, and thus is harmonic.

I will first introduce some tools that are required for the analysis of this problem, and are also useful in the broader field of analysis. The main analytical tool used in this study is the Fourier transform. The Fourier transform is then used to derive the Poisson kernel, which provides an explicit harmonic expression for the extension of the boundary data in the form of a Poisson integral.

The thesis presents two main results. The first of these proves that the Poisson integral can be used as an extension of a boundary value function when it belongs to the $L^p$ space. An extension of this result with stricter regularity assumptions on the function, which leads to a stronger notion of convergence, is also presented. The second main result shows that a similar result also holds for finite Borel measures. This allows an even wider variety of data to be represented at the boundary of the half-space.

These results provide an explicit solution to a boundary value problem for Laplace's equation, but also show how harmonic extensions behave under different regularity assumptions. These results have applications in a wide range of mathematical analysis, especially in partial differential equations, and in physics in the form of potential theory. In the end, I will highlight some of the possible ways in which these results can be extended to cover some more generalized or slightly modified settings.
\end{foreignlanguage}


\vfill

\noindent\textbf{Keywords:} Harmonic analysis, Partial differential equations, Fourier analysis % Avainsanat / Nyckelord

\clearpage

%
% TABLE OF CONTENTS
%
\cleardoublepage
\tableofcontents


%
% MAIN TEXT STARTS HERE
%

% The \pagenumbering{arabic} command must be immediately after the first chapter title.
% Do not repeat it for other chapters.
\chapter{Introduction}
\pagenumbering{arabic}
In this thesis, I will study results in the field of harmonic and Fourier analysis. I will be concentrating on results that prove that Poisson integrals in the case of $L^p$ spaces and Poisson-Stieltjes integrals in the case of finite Borel measures are a way to extend a function describing some boundary value at the boundary of a half space $\RR^n\times\{0\}$, to the whole of the half space in a way that it satisfies the Laplace equation $\Delta u=\sum_i \frac{\partial^2 u}{\partial x_i^2}=0$, while converging towards the function at the boundary. This work will follow mostly the results given in the \textit{Introduction to Fourier Analysis on Euclidean Spaces} by E. M. Stein and G. Weiss \cite{stein_weiss}.

The problem that is investigated in this thesis is a special case of the Dirichlet problem. The problem poses a question: can we find a solution for a partial differential equation (PDE), such that it will converge to some specified function at the boundary of the space? For the purpose of this thesis, the PDE that will be considered is the Laplace equation, and the space will be the half space $\RR^n\times\RR_+$.

This question relates to many areas of mathematical analysis and is deeply connected to mathematical physics. In mathematics, these types of results are frequently used in the analysis of PDEs to extend certain weaker solution concepts to boundary value problems. This allows embedding less regular initial value data, which may not satisfy the boundary condition in the classical sense of functions. These are connected to physical problems through potential theory, where, for example, gravitational and electrostatic potentials can be described by harmonic functions.

In the rest of this thesis, I will first introduce some of the essential tools that are used in the construction of the Poisson kernel in Chapter \ref{ch:2}. This will also include showing some properties of the Poisson kernel that will be needed later. Chapter \ref{ch:3} will concentrate on the main results of this thesis. First, I will show that the Poisson integral extends the boundary value function in $L^p$ to the half space. The second part will extend this result to also include finite Borel measures.

\chapter{Poisson Formula for Harmonic Functions in $\RR_+^{n+1}$}\label{ch:2}
In order to define the Poisson kernel

%\begin{definition}[Summability to $L$]\label{def:mean_of_int}
%    A function $f$ is said to be $M$-summable to $L$, if $\lim_{\epsilon\rightarrow0}M_{\epsilon,\Phi}(x)=L$, where
%
%    \begin{equation*}
%        M_{\epsilon,\Phi}(f)=M_{\epsilon}(f)=\int_{\RR^n}\Phi(\epsilon x)f(x)dx
%    \end{equation*}

%    \noindent where $\Phi\in C_0$ and $\Phi(0)=1$. Then $M_\epsilon(f)$ is called the $\Phi$-mean of the integral.
%\end{definition}

\clearpage
\chapter{Harmonic functions in $\RR_+^{n+1}$}\label{ch:3}

\begin{theorem}
    If $\mu\in M(E_+^{n+1})$ and $u(x,y)=\int_{E_n}P(x-t,y)d\mu(t)$ is the Poisson-Stiltjes integral of $\mu$ then $u$ is harmonic in $E_+^{n+1}$ and 

    \begin{equation}
        ||u(\cdot,y)||_1=\int_{E_n}|u(x,y)|dx\leq||\mu||
    \end{equation}

    \noindent where $||\mu||$ is the total variation, $|\mu|(E_n)$, of $\mu$. Furthermore,

    \begin{equation*}
        \lim_{y\rightarrow0}\int_{E_n}u(x,y)\varphi(x)dx=\int_{E_n}\varphi(x)d\mu(x)
    \end{equation*}

    \noindent for all $\varphi\in C_0$; that is, $u(\cdot,y)$ converges to $\mu$ in the weak* topology.
\end{theorem}

\clearpage
\chapter{Conclusion}

This concludes the results in the upper half space with tangential convergence to a function or measure. Together, theorems \ref{thm:21a} \& \ref{thm:2.3} show that there exists an extension of a function $L^p(\RR^n)$ or a Borel measure $\mathcal M(\RR^n)$ to the half-space $\RR^{n+1}_+$ that solves the Laplace equation and thus is harmonic. Theorem \ref{thm:21b} also shows that the result can be strengthened by having stronger assumptions about the regularity of the boundary data.

These results are only a first step to the results that could be further generalized. One possibility is to consider other paths to the boundary that do not necessarily approach it tangentially. Stein and Weiss are also considering this scenario in their book \cite{stein_weiss}. With a nontangential convergence path, the analysis becomes more complicated and requires some tools that were not used for the theorems introduced in this thesis. In particular, this would require the use of maximal functions to obtain certain bounds to establish the result in the non-tangential case. Other possible extensions could include changing the domain in which the problem is set. Also, the differential equation could be changed from the Laplace equation to some other partial differential equation.

%
% BIBLIOGRAPHY
%
% The [heading=bibintoc] option adds the bibliography to table of contents
\printbibliography[heading=bibintoc]




%
% APPENDICES (if necessary)
%
%\appendix

%\chapter{Some more information}

%If you need to have appendices in your thesis,
%you can write them as usual chapters after the \verb|\appendix| command.
%Appendices are numbered alphabetically.

\end{document}
